\chapter{单变量函数的微分学}
\section{导数}
\subsection{导数的定义}
\begin{definition}
    设$f(x)$在$x_0$的邻域内有定义，若$\lim\limits_{x\to x_0}\dfrac{f(x)-f(x_0)}{x-x_0}$存在且有限，那么称它为$f(x)$在$x_0$的\overtext{导数}{derivative}，记作：
    \[f'(x_0)\quad\text{或}\quad\eval{\dv{f}{x}}_{x_0}\quad\text{或}\quad\eval{\dv{y}{x}}_{x_0}\]
    并称$f(x)$在$x_0$处\overtext{可导}{differentiable}.
    \begin{itemize}
        \item 若有$\lim\limits_{x\to x^+_0}\dfrac{f(x)-f(x_0)}{x-x_0} = \lim\limits_{\Delta x \to 0^+}\dfrac{f(x_0+\Delta x)-f(x_0)}{\Delta x} = f_+'(x_0)$存在且有限，称之为\textbf{右导数}.
        \item 若有$\lim\limits_{x\to x^-_0}\dfrac{f(x)-f(x_0)}{x-x_0} = \lim\limits_{\Delta x \to 0^-}\dfrac{f(x_0+\Delta x)-f(x_0)}{\Delta x} = f_-'(x_0)$存在且有限，称之为\textbf{右导数}.
    \end{itemize}
    若$f$在$I$上任意一点都可导，则称$f$在$I$上可导.
\end{definition}

\begin{theorem}
    $f(x)$在$x_0$可导，当且仅当$f_\pm'(x_0)$存在且i相等.
\end{theorem}

\begin{theorem}
    $f(x)$在$x_0$可导，则$f(x)$在$x_0$连续.
\end{theorem}
\begin{proof}
    \[\lim_{x \to x_0}(f(x)-f(x_0)) = \lim_{x \to x_0}\left[\frac{(f(x)-f(x_0))}{x-x_0}\right](x-x_0) = f'(x)\cdot 0 = 0\]
\end{proof}


\subsection{导数的四则运算}
\begin{theorem}
    设$f(x)$在$x$可导，则：
    \begin{align*}
        (f(x)\pm g(x))' &= f'(x) \pm g'(x) \\
        (f(x)g(x))' &= f'(x)g(x) + f(x)g'(x) \\
        \left(\frac{f(x)}{g(x)}\right)' &= \frac{f'(x)g(x)-f(x)g'(x)}{g^2(x)}\,(g(x)\neq 0)
    \end{align*}
\end{theorem}
以乘法的证明为例：
\begin{proof}
    设$x_0 \in I$：
    \begin{align*}
        \frac{f(x)g(x)-f(x_0)g(x_0)}{x-x_0} &= \frac{f(x)g(x)-f(x_0)g(x)+f(x_0)g(x)-f(x_0)g(x_0)}{x-x_0}\\
        &= \frac{f(x)-f(x_0)}{x-x_0}g(x)+\frac{g(x)-g(x_0)}{x-x_0}f(x_0) \\
        &= f'(x_0)g(x) + f(x_0)g'(x_0)
    \end{align*}
    由于$g(x)$可导，其必然连续，于是$x\to x_0$时有$g(x)\to g(x_0)$，代入上式得证.
\end{proof}


\subsection{复合函数的求导法则}
\begin{theorem}
    设$y=g(x)$定义在区间$I$上，$z=f(y)$定义在区间$J$上，且$g(I) \subset J$，如果$g(x)$在$x\in I$处可导，$f(y)$在$y=g(x)$处可导，那么$f \circ g$在$x$处可导，且有：
    \[(f \circ g)'(x)=f'(g(x))g'(x) \quad\text{或}\quad \dv{z}{x}=\dv{z}{y}\dv{y}{x}\]
\end{theorem}

\begin{proof}
    取$x_0 \in I$记$y_0 = g(x_0)$，先考虑$x_0$和$g(x_0)$不在所在定义域的两端点的情况，其余情况修改部分定义易证.注意到：
    \begin{align*}
        \dv{z}{x} =& \lim_{x\to x_0}\frac{f(g(x))-f(g(x_0))}{x-x_0}\\
        &= \lim_{x\to x_0}\frac{f(g(x))-f(g(x_0))}{g(x)-g(x_0)}\frac{g(x)-g(x_0)}{x-x_0}
    \end{align*}
    在上式中，若是$g(x)-g(x_0)=0$将导致上式中用于辅助的中间项出现问题，这是构造过程中导致的问题，而不是原式所具有的性质，所以要消除这个影响，尝试构建函数$h(x)$：
    \[h(x)=\begin{cases}
        \dfrac{f(y)-f(y_0)}{y-y_0},\quad y\neq y_0\\
        f'(y),\quad y=y_0
    \end{cases}\]
    若是使用$h(x)$重新构造展开式，则$g(x)-g(x_0)=0$将能够被$(f\circ g)'(x)$接受，避免被排出定义域中：
    \[\dv{z}{x} = \lim_{x\to x_0}h(g(x))\frac{g(x)-g(x_0)}{x-x_0}\]
    由于$h(x)$连续，所以有$\lim\limits_{x\to x_0}h(g(x)) = h\left(\lim\limits_{x\to x_0}g(x)\right) = h(g(x_0))$，故而上式可以进一步写作：
    \[\dv{z}{x} = f'(g(x))g'(x)\]
\end{proof}


\subsection{反函数的求导法则}
\begin{theorem}
    \[\left(f^{-1}(y)\right)' = \dv{x}{y} = \left(\dv{y}{z}\right)^{-1} = \frac{1}{f'(x)}\]
\end{theorem}

\begin{proof}
    \[\frac{f^{-1}(y)-f^{-1}(y_0)}{y-y_0} = \frac{x-x_0}{y-y_0} = \left(\frac{y-y_0}{x-x_0}\right)^{-1} \to \frac{1}{f'(x)}\]
\end{proof}

利用反函数的求导，可以方便地证明各种初等函数尤其是对数、指数、反三角函数的导数.

\begin{question}
    求$y = x \ee^x$的反函数的导数.
\end{question}

\begin{solution}
    代入$x=f(y)$得$y = f(y)\ee^{f(y)}$，两边对$y$求导：
    \begin{align*}
        1 &= f'(y)\ee^{f(y)} + f(y)f'(y)\ee^{f(y)}\\
        f'(y) &= \frac{1}{\ee^{f(y)} + f(y)\ee^{f(y)}} \\
        f'(y) &= \frac{1}{\ee^x + x\ee^x},\quad\text{其中$x$满足$y=x\ee^x$}
    \end{align*}
\end{solution}

在完全求不出反函数的时候，直接用一个抽象函数$x=f(y)$代表反函数并且代入到原函数表达式的$x$中，同时两边对$y$求导，最后再把抽象的$f'(y)$与$f(y)$分离，最后再把$f(y)=x$代回，最终获得表达式.